\documentclass[10pt, letterpaper]{article}

% Packages:
\usepackage[
    ignoreheadfoot, % set margins without considering header and footer
    top=1.25 cm, % separation between body and page edge from the top
    bottom=1.25 cm, % separation between body and page edge from the bottom
    left=1.25 cm, % separation between body and page edge from the left
    right=1.25 cm, % separation between body and page edge from the right
    footskip=1.0 cm, % separation between body and footer
    % showframe % for debugging 
]{geometry} % for adjusting page geometry
\usepackage{titlesec} % for customizing section titles
\usepackage{tabularx} % for making tables with fixed width columns
\usepackage{array} % tabularx requires this
\usepackage[dvipsnames]{xcolor} % for coloring text
\definecolor{primaryColor}{RGB}{0, 0, 0} % define primary color
\usepackage{enumitem} % for customizing lists
\usepackage{fontawesome5} % for using icons
\usepackage{amsmath} % for math

\usepackage[
    pdftitle={Harry's CV},
    pdfauthor={Harry},
    pdfcreator={LaTeX with RenderCV},
    colorlinks=true,
    urlcolor=primaryColor
]{hyperref} % for links, metadata and bookmarks
\usepackage[pscoord]{eso-pic} % for floating text on the page
\usepackage{calc} % for calculating lengths
\usepackage{bookmark} % for bookmarks
\usepackage{lastpage} % for getting the total number of pages
\usepackage{changepage} % for one column entries (adjustwidth environment)
\usepackage{paracol} % for two and three column entries
\usepackage{ifthen} % for conditional statements
\usepackage{needspace} % for avoiding page brake right after the section title
\usepackage{iftex} % check if engine is pdflatex, xetex or luatex

% Ensure that generate pdf is machine readable/ATS parsable:
\ifPDFTeX
    \input{glyphtounicode}
    \pdfgentounicode=1
    \usepackage[T1]{fontenc}
    \usepackage[utf8]{inputenc}
    \usepackage{lmodern}
\fi

\usepackage{charter}

% Some settings:
\raggedright
\AtBeginEnvironment{adjustwidth}{\partopsep0pt} % remove space before adjustwidth environment
\pagestyle{empty} % no header or footer
\setcounter{secnumdepth}{0} % no section numbering
\setlength{\parindent}{0pt} % no indentation
\setlength{\topskip}{0pt} % no top skip
\setlength{\columnsep}{0.15cm} % set column separation
\pagenumbering{gobble} % no page numbering

\titleformat{\section}{\needspace{4\baselineskip}\bfseries\large}{}{0pt}{}[\vspace{1pt}\titlerule]

\titlespacing{\section}{
    % left space:
    -1pt
}{
    % top space:
    0.15 cm
}{
    % bottom space:
    0.1 cm
} % section title spacing

\renewcommand\labelitemi{$\vcenter{\hbox{\small$\bullet$}}$} % custom bullet points
\newenvironment{highlights}{
    \begin{itemize}[
        topsep=0.05 cm,
        parsep=0.05 cm,
        partopsep=0pt,
        itemsep=0pt,
        leftmargin=0 cm + 10pt
    ]
}{
    \end{itemize}
} % new environment for highlights

\newenvironment{highlightsforbulletentries}{
    \begin{itemize}[
        topsep=0.05 cm,
        parsep=0.05 cm,
        partopsep=0pt,
        itemsep=0pt,
        leftmargin=10pt
    ]
}{
    \end{itemize}
} % new environment for highlights for bullet entries

\newenvironment{onecolentry}{
    \begin{adjustwidth}{
        0 cm + 0.00001 cm
    }{
        0 cm + 0.00001 cm
    }
}{
    \end{adjustwidth}
} % new environment for one column entries

\newenvironment{twocolentry}[2][]{
    \onecolentry
    \def\secondColumn{#2}
    \setcolumnwidth{\fill, 4.5 cm}
    \begin{paracol}{2}
}{
    \switchcolumn \raggedleft \secondColumn
    \end{paracol}
    \endonecolentry
} % new environment for two column entries

\newenvironment{threecolentry}[3][]{
    \onecolentry
    \def\thirdColumn{#3}
    \setcolumnwidth{, \fill, 4.5 cm}
    \begin{paracol}{3}
    {\raggedright #2} \switchcolumn
}{
    \switchcolumn \raggedleft \thirdColumn
    \end{paracol}
    \endonecolentry
} % new environment for three column entries

\newenvironment{header}{
    \setlength{\topsep}{0pt}\par\kern\topsep\centering\linespread{1.1}
}{
    \par\kern\topsep
} % new environment for the header

\newcommand{\placelastupdatedtext}{% \placetextbox{<horizontal pos>}{<vertical pos>}{<stuff>}
  \AddToShipoutPictureFG*{% Add <stuff> to current page foreground
    \put(
        \LenToUnit{\paperwidth-2 cm-0 cm+0.05cm},
        \LenToUnit{\paperheight-1.0 cm}
    ){\vtop{{\null}\makebox[0pt][c]{
        \small\color{gray}\textit{Last updated in September 2024}\hspace{\widthof{Last updated in September 2024}}
    }}}%
  }%
}%

% save the original href command in a new command:
\let\hrefWithoutArrow\href

\begin{document}
    \newcommand{\AND}{\unskip
        \cleaders\copy\ANDbox\hskip\wd\ANDbox
        \ignorespaces
    }
    \newsavebox\ANDbox
    \sbox\ANDbox{$|$}

    \begin{header}
        \fontsize{25 pt}{25 pt}\selectfont Harry Zhang

        \vspace{5 pt}

        \normalsize
        \mbox{Toronto, Ontario}%
        \kern 5.0 pt%
        \AND%
        \kern 5.0 pt%
        \mbox{\hrefWithoutArrow{mailto:harryh.zhang@mail.utoronto.ca}{harryh.zhang@mail.utoronto.ca}}%
        \kern 5.0 pt%
        \AND%
        \kern 5.0 pt%
        \mbox{\hrefWithoutArrow{tel:+1-647-213-3096}{+1-647-213-3096}}%
        \kern 5.0 pt%
        \AND%
        \kern 5.0 pt%
        \mbox{\hrefWithoutArrow{https://www.linkedin.com/in/harryhnzhang/}{www.linkedin.com/in/harryhnzhang/}}%
        
    \end{header}

    \vspace{0.15 cm}

    \section{Education}

        \begin{twocolentry}{
            Sept 2024 -- Apr 2029
        }
            \textbf{University of Toronto}, BASc in Computer Engineering, PEY Co-op\end{twocolentry}

        \vspace{0.05 cm}
        \begin{onecolentry}
         \begin{highlights}
                \item \textbf{cGPA:} 3.65/4.00
                \item \textbf{Coursework:} Programming Fundamentals, Computer Organizations, Electronics, Digital Systems, Signals and Systems, Linear Algebra, Calculus I, II \& III 
            \end{highlights}
        \end{onecolentry}

    \section{Experience}

        \begin{twocolentry}{Jan 2025 -- Apr 2025}
            \textbf{Undergraduate Research Assistant, Data-Capture Systems Testing \& Validation}
            
            \text{Wellness and Health Enhancement Engineering Lab (WHEEL), University of Toronto}
        \end{twocolentry}

        \vspace{0.05 cm}
        \begin{onecolentry}
            \begin{highlights}
                \item Supervised by Dr. Enid Montague, collaborated with UofT's Wellness and Health Enhancement Engineering (WHEEL) Lab to redesign a usability testing lab for studying physician-technology interactions
                \item Tested and validated Livestreamer's prototype data capture system by executing multimodal performance tests, achieved strong multimodal coverage (0.91).
                \item Coordinated quality assurance of 3D lab layout models using Fusion 360 and Blender, improved communication with clients and enabled accurate assessment of feasibility.
            \end{highlights}
        \end{onecolentry}
        
    \section{Projects}

         \begin{twocolentry}{May 2026 -- Aug 2026}
            \textbf{Pose-Based Human Action Recognition}\end{twocolentry}
        \begin{onecolentry}{APS360 (Applied Fundamentals of Deep Learning), University of Toronto}\end{onecolentry}

             \vspace{0.05 cm}
            \begin{onecolentry}
                \begin{highlights}
                \item Built a \textbf{recurrent neural network} in \textbf{PyTorch} to classify 30-frame body-pose sequences (17 COCO keypoints) into 20 action classes on the MPOSE2021 dataset, modeling spatio-temporal joint dynamics.
                \item Engineered a keypoint preprocessing pipeline that hip-centers and torso-scales each frame and masks low-confidence joints, with a stratified, subject-independent 70/15/15 split to prevent data leakage.
                \item Benchmarked the RNN against an averaged-pose \textbf{MLP} baseline using weighted cross-entropy, dropout, and early stopping to address class imbalance and overfitting.
            \end{highlights}
        \end{onecolentry}

        \vspace{0.15 cm}

         \begin{twocolentry}{Jan 2026 -- Apr 2026}
            \textbf{Geographic Information System (GIS) for City of Toronto}\end{twocolentry}
        \begin{onecolentry}{ECE297 Software Design Project, University of Toronto}\end{onecolentry}

             \vspace{0.05 cm}
            \begin{onecolentry}
                \begin{highlights}
                \item Applied \textbf{Git/GitHub} and \textbf{VS Code debugger} to manage iterative development, debug logic/memory issues, and validate performance, ensured stable map loading and efficient query execution
                \item Developed modular C++ solutions using \textbf{STL containers and algorithms} to process OSM data; wrote clean, maintainable code and reduced runtime errors during large-scale dataset testing.
                \item Wrote clear technical documentation and researched OpenStreetMap data structures by studying the API and analyzing map data formats, reduced implementation errors and kept the team aligned during development.
            \end{highlights}
        \end{onecolentry}

        \vspace{0.15 cm}
        
         \begin{twocolentry}{Nov 2025}
            \textbf{Connect8 Puzzle Game on FPGA}\end{twocolentry}
        \begin{onecolentry}{ECE241 (Digital Systems) Course Project, University of Toronto}\end{onecolentry}

             \vspace{0.05 cm}
            \begin{onecolentry}
                \begin{highlights}
                \item Built a full hardware puzzle game in Verilog on the DE1-SoC by implementing block movement, collision detection, and an 8x8 line-clear engine, produced a stable core game loop that ran entirely on hardware.
                \item Integrated a PS/2 keyboard controller, \textbf{VGA renderer}, and a 16-bit \textbf{LFSR block generator} to support real-time input and display, gave the game smooth controls and a consistent 60 FPS output.                
                \item Debugged critical rotation and collision handling errors and revised code based on testbench validation
            \end{highlights}
        \end{onecolentry}

    \section{Skills}
        
        \begin{onecolentry}
            \textbf{Programming}: C, C++, Python, RISC-V Assembly (NIOS V), MATLAB, Verilog
        \end{onecolentry}
        
        \vspace{0.1 cm}

        \begin{onecolentry}
            \textbf{Technologies}: VS Code, Git/Github, AutoCAD
        \end{onecolentry}

    \section{Honors \& Awards}
    
    \begin{twocolentry}{Fall 2024, Winter \& Fall 2025}
            \textbf{Dean's List}\end{twocolentry}
            
    \vspace{0.05 cm}
    \begin{twocolentry}{\$10,000}
            \textbf{Faculty Of Applied Science And Engineering Admission Scholarship}\end{twocolentry}
            
    \vspace{0.05 cm}
    \begin{twocolentry}{\$2,500}
            \textbf{The Edward S. Rogers Sr. Department Admission Scholarship}\end{twocolentry}

\end{document}